\chapter{The legal issue: Clipper's Still Kicking}

\section{Introduction}
Although the Clipper chip controversy seems to have died down, the
Clipper initiative is not dead. The U.S. government's attempt to
legislate an encryption technology that gives it a backdoor to monitor
all our transmissions still lives. That backdoor is Clipper.\\

The argument is as old as encryption. On one side, you have those in
authority who, for the good of the state, want to be able to read or
listen to anyone's communications. On the other, you have civil
libertarians who insist that everyone has a right to privacy against
unreasonable searches.\\

This conflict exploded into the Internet's consciousness in 1994 when
Vice President Al Gore and FBI director Louis Freeh pushed for the use
of the Clipper chip and the passage of the Digital Telephony bill. Most
of the public furor was directed at the chip, but Clipper and the
Digital Telephony bill work hand in glove with each other.

\section{The Clipper Controversy}
Before examining the current state of Clipper, let's re-examine the
root of the controversy. Clipper is a computer chip that encodes
digital voice and data communications using an encryption algorithm
called Skipjack, which is a product of the National Security Agency
(NSA). The chip is meant for use in telephones, fax machines, and
modems.\\

What makes Clipper different from other encryption devices, and what
has drawn the ire of many citizens, is that the chip has a trap door
that government investigators can open to wiretap Clipper-equipped
devices. The door's key, which is unique to each chip, has two parts.
In the government's plan, the two halves of the key would be held
separately in escrow by the Treasury and Commerce Departments. To use
the keys, federal agents would need to obtain a warrant.\\

It sounds safe enough, but more than 45,000 Internet users disagreed and
signed an electronic petition against Clipper. Many people fear the
government could misuse Clipper to harm innocent people and businesses.\\

The FBI and NSA argued that the government was not about to become the
Peeping Tom of the online world and that the security holes were needed
for criminal investigations. Civil libertarians on the Internet replied
that Clipper and the concepts behind it were unacceptable. They argued
that even if the security holes were used for acceptable purposes,
there was no guarantee that future administrations might not use this
against them or their children. In particular, they objected to the
keys being held by the government. But the federal government, in
essence, shrugged and continued to push for Clipper acceptance in
telecom and datacom businesses.\\

Clipper was not home free yet. Against a mounting storm of protest, the
executive branch found Clipper under siege from all sides. This led to
such odd alliances as the Electronic Freedom Foundation (EFF) working
with conservative Republican senators against the Clipper initiative.\\

All this protest might have gone for naught until AT\&T Bell Labs
announced that one of its engineers had shown that Skipjack could be
broken. Against a rising storm of protest, Vice President Gore wrote a
letter on July 20, 1994 to House of Representatives member Maria
Cantwell (Wash.) in which he said he believed there must be further
public comment on Clipper.\\

Cantwell seized on this as a face-saving maneuver signaling the end of
the executive branch's support for Clipper and announced victory for
the anti-Clipper forces. She soon was seconded in this interpretation
by the EFF, and the Clipper controversy sank from sight.

\section{A Chip with Nine Lives}
But the Clipper issue has not died. The NSA continues to work on an
encryption algorithm that will be proof against all comers--except
federal agents armed with the decryption key. Clipper itself, flaws and
all, is still on the books as a voluntary Federal Information
Processing Standard (FIPS) for government telecommunications.\\

If you want to export a telephone device today, you still have no
choice but to ship Clipper-equipped products. Of course, finding a
foreign customer that would buy such a machine might prove difficult.\\

Clipper, it should also be noted, was never designed to work with local
area networks or the Internet. The NSA and National Institute of
Science and Technology (NIST) are working to develop software
encryption systems with key-escrow backdoors that would work on
networks.\\

In short, has anything really changed? It doesn't look like it. No one
(except employees in federal agencies that have adopted Clipper
equipment) is required to use a key-escrow security system.
Nevertheless, looks, as we shall see, can be deceiving.\\

But, you might ask, does it really matter? After all, we have
public-key encryption systems using the Rivest-Shamir-Adleman (RSA)
encryption algorithm. If you don't trust RSA--and there's reason to
believe it might have a hole in it--there are other ways to prevent
someone from reading your e-mail. For instance, Philip Zimmermann's
popular Pretty Good Privacy (PGP) program uses RSA, International Data
Encryption Algorithm (IDEA), and Message Digest Algorithm \#5 (MD5) to
triply ensure that your messages remain private. Zimmermann also is
working on a voice PGP version to give telephone customers an
alternative to Clipper.\\

PGP, however, like other encryption schemes, cannot be exported legally
from the United States. As a result, Zimmermann is under investigation
by the U.S. Customs Service for possible violations of the
International Traffic in Arms Regulations (ITAR). It is conceivable
that because it's clearly impossible to keep information from flowing
freely on the Internet, there will be an attempt to regulate PGP. A
requirement to register keys with the government seemingly would fit
perfectly into Clipper's key-escrow framework.\\

Does this sound unlikely? Consider then, that Congress passed the
FBI-backed Digital Telephony bill in late October last year. This
legislation requires that common carriers, local telephone companies,
and long-distance services must add openings for federal wiretaps to
their phone systems. The bill also authorizes government agencies to
access billing records. At least police and federal agencies will need
search warrants to tap your digital lines and dig through your bills.

\section{Uncle Sam is Listening}
The telephone industry, after fighting this act tooth and nail, is
slowly getting used to the idea. Its problem with the bill is not so
much privacy considerations as a well-grounded fear that retrofitting
older digital equipment and new digital phone systems will cost it a
fortune in unnecessary expenses. The government will be supplying \$500
million for retrofitting existing systems, but the telephone companies
fear this will only be the first down payment on a bill that could
swell to the billions of dollars.\\

The telephone companies have some small protection from a never-ending
spiral of wiretapping bills. The legislation only requires that they
pay `reasonably achievable ` costs. The Federal Communications
Commission (FCC) will decide, on a case by case basis, what costs fall
into this category. If the FCC judges the costs unreasonable, a
telephone company could petition the attorney general for funds to
accomplish the operation. If the attorney general doesn't reimburse the
company, the phone company could refuse to comply, remaining on the
right side of the law.

\section{Getting Clipped}
The telephone industry isn't the only potential target of the Digital
Telephony bill. Any public communications network that carries both
voice and data also is affected. For example, if TV coaxial networks
became avenues for two-way communications, they would be subject to the
bill's regulations. Whether the Internet, with its relatively primitive
voice-communications facilities, should be immune to this bill's
regulations hasn't been decided.\\

There is a catch, however. If you `re calling your Internet provider
via a long-distance line or you `re using an Integrated Services
Digital Network (ISDN), your Internet communications would be
susceptible to Digital Telephony-authorized taps.\\

Under the bill, telecom carriers not only have to ensure that the
government can tap their communication lines, they have to aid the
government in its snooping. Telephone companies are required to be able
to find a suspect's numbers and report where his or her phone calls are
originating and where they're going.\\

The bill could have been worse. One version would have given the
attorney general the right to stop any technological advance that could
interfere with the government's ability to make wiretaps. Another early
model of the bill would not have required search warrants to engage in
wiretapping.\\

\section{An Orwellian Future?}
Even from a civil libertarian point of view, there are a few things
right about the Digital Telephony bill. Most notable is that a police
officer or federal agent must obtain a court order to access your
e-mail. In the past, all that was required was a subpoena, and
generally speaking, subpoenas are easier and faster for an officer to
obtain.\\

What does it all mean for you as an Internet user? In the short run,
probably nothing. Look down the road a few years and it's another
story. In the future, all telecom systems will be digital and, thus,
tappable. It's doubtful that federal law enforcement agencies will put
up with being able to easily tap communications without having the
ability to decode the messages.

